\section{Introduction}

Reproducible image-analysis software can be wrong in a particularly durable way. A fixed environment, deterministic seeds, complete hashes, and byte-identical outputs show that a computation can be repeated; they do not show that a named feature is responsive to image content. Optional computer-vision modules make this distinction consequential: a missing namespace can be caught, replaced with a constant, and propagated through thousands of images while every conventional integrity check still passes.

This risk is not hypothetical for the visual-communication and image-analysis community. Recent JVCIR and related venues increasingly combine handcrafted visual descriptors---edge statistics, colour moments, saliency maps, layout metrics---with computational aesthetics, complexity prediction, and design-quality assessment \citep{guo2026dual,yang2024selfsupervised,yan2024hybrid}. These pipelines often depend on optional libraries, and their published benchmarks are expected to be reproducible. Yet reproducibility guarantees bind only the computation that was actually executed, not the computation that was intended.

This paper studies one such failure in a released visual-feature pipeline. The extractor attempted to create an OpenCV fine-grained saliency object. In environments containing core OpenCV but not its optional contrib namespace, the call failed and a broad exception handler returned a constant 0.5 map. Six downstream fields then became constant across 16,425 Pillow-rendered images and 1,250 images from a separate OpenCV renderer. The affected fields carried 15\% of the absolute weight in a legacy order formula and 45\% in a legacy hierarchy formula. Hashes faithfully reproduced the defect.

The failure motivates three research questions:
\begin{enumerate}
    \item \textbf{RQ1:} Which development-time semantic-health checks expose inactive weighted features and rank-duplicate baselines that finite-value and hash checks miss?
    \item \textbf{RQ2:} How much can repairing a silent feature fallback change held-out conclusions when target scaling, feature scaling, and baseline choice are confined to development data?
    \item \textbf{RQ3:} Do repaired signals remain nondegenerate under a separately implemented renderer, and what do complete mechanism--target crossing tables reveal about attribution and scope?
\end{enumerate}

The five weighted composites in the case study are historical software artefacts named complexity, order, harmony, intensity, and hierarchy. They are not externally sourced perceptual scales; their value here is diagnostic---they make the propagation of a feature-layer defect observable at formula and conclusion levels. The rendering-parameter targets are likewise controlled engineering variables, not substitutes for perceived complexity, aesthetics, or preference.

The paper makes four contributions. First, it provides reusable semantic-health gates for non-finite or constant weighted inputs, retained absolute formula weight, exact ties and clipping, and rank-equivalent baseline candidates. These gates complement existing metamorphic and property-based testing approaches for scientific software by targeting the specific failure mode of semantically inactive but numerically complete feature columns. Second, it replaces the optional-module path with an in-repository spectral-residual implementation and changes missing, incomplete, or degenerate computation from silent substitution to explicit failure. Third, it reports a pre/post failure analysis under a leakage-controlled protocol: target and feature constants are fitted on development data only, elementary baselines are selected on development data and locked before held-out evaluation, all 25 proxy--target associations are disclosed, and uncertainty resamples rendering seeds as clusters. Fourth, it repeats semantic and mechanism checks in a second internally controlled renderer and releases executable code, images, manifests, tests, environment locks, and artifact hashes.

The intended contribution is therefore an image-analysis software result relevant to JVCIR's scope in visual-feature computation and benchmark reliability: deterministic reproducibility needs semantic invariants that test whether intermediate representations remain meaningfully connected to the image. The empirical results quantify one failure and repair across two controlled rendering systems. They do not establish natural-image generalisation or human perceptual validity.

